% MV-HGS-SP main algorithm flow. Transparent line art only.
% The sequence is bound to route_pool_sp.py: three independent HGS views,
% full-model archive review, a time-limited MIP route-pool recombination,
% full-model verification, and monotone acceptance.
\begin{figure}[H]
\centering
\begin{tikzpicture}[
    every node/.style={font=\fontsize{8.2}{9.6}\selectfont},
    main/.style={rectangle,draw,thin,text width=6.5cm,minimum height=7.4mm,
        align=center,inner xsep=3pt,inner ysep=2pt},
    view/.style={rectangle,draw,thin,text width=3.35cm,minimum height=10.5mm,
        align=center,inner xsep=2pt,inner ysep=1.8pt},
    review/.style={rectangle,draw,thin,text width=3.35cm,minimum height=11.6mm,
        align=center,inner xsep=2pt,inner ysep=1.8pt},
    decision/.style={diamond,draw,thin,aspect=3.35,text width=4.15cm,
        align=center,inner sep=1.0pt},
    result/.style={rectangle,draw,thin,text width=3.45cm,minimum height=7.2mm,
        align=center,inner xsep=2pt,inner ysep=1.8pt},
    edge/.style={-{Stealth[length=2.0mm,width=1.35mm]},line width=0.46pt,
        line cap=round,line join=miter},
    stem/.style={line width=0.46pt,line cap=round,line join=miter},
    labeltext/.style={font=\fontsize{8.0}{9.2}\selectfont,inner sep=1.1pt}
]
\node[main] (input) at (0,0)
    {读取算例、时空参数、有限车队与统一评价设置};
\node[main] (fallback) at (0,-1.05)
    {构造共同可行后备解};

\coordinate (split) at (0,-1.72);
\node[labeltext] at (0,-1.98)
    {三个机制视角独立搜索（相同种子、独立种群）};
\node[view] (hgsf) at (-4.15,-2.78)
    {HGS-F\\燃油成本视角\\固定5000次迭代};
\node[view] (hgse) at (0,-2.78)
    {HGS-E\\简化电动视角\\固定5000次迭代};
\node[view] (hgsm) at (4.15,-2.78)
    {HGS-M\\机制感知视角\\固定5000次迭代};

\node[review] (reviewf) at (-4.15,-4.20)
    {完整模型复核\\24个档案候选\\及后备、代理最优解};
\node[review] (reviewe) at (0,-4.20)
    {完整模型复核\\24个档案候选\\及后备、代理最优解};
\node[review] (reviewm) at (4.15,-4.20)
    {完整模型复核\\24个档案候选\\及后备、代理最优解};

\coordinate (merge) at (0,-5.03);
\node[main,text width=9.15cm] (archive) at (0,-5.68)
    {每视角保留8个完整成本最低的父方案；汇集全部复核成功候选的路线};
\node[main] (pool) at (0,-6.78)
    {路线合并去重，建立完整成本路线池};
\node[main] (mip) at (0,-7.88)
    {HiGHS限时MIP路线重组（5~s）；记录可行解、下界、间隙与状态};
\node[main] (verify) at (0,-8.98)
    {重组方案完整补全、客户恰好覆盖检查与独立核验};
\node[decision] (better) at (0,-10.30)
    {重组方案严格优于最好父方案};

\node[result] (accept) at (-2.55,-11.98)
    {接受重组方案};
\node[result] (retain) at (2.55,-11.98)
    {保留最好父方案};
\coordinate (branch) at (0,-11.30);
\coordinate (outmerge) at (0,-12.70);
\node[main] (output) at (0,-13.37)
    {输出总体最好可行解};

\draw[edge] (input) -- (fallback);
\draw[stem] (fallback.south) -- (split);
\draw[edge] (split) -| (hgsf.north);
\draw[edge] (split) -- (hgse.north);
\draw[edge] (split) -| (hgsm.north);
\draw[edge] (hgsf) -- (reviewf);
\draw[edge] (hgse) -- (reviewe);
\draw[edge] (hgsm) -- (reviewm);
\draw[stem] (reviewf.south) |- (merge);
\draw[stem] (reviewe.south) -- (merge);
\draw[stem] (reviewm.south) |- (merge);
\draw[edge] (merge) -- (archive);
\draw[edge] (archive) -- (pool);
\draw[edge] (pool) -- (mip);
\draw[edge] (mip) -- (verify);
\draw[edge] (verify) -- (better);
\draw[stem] (better.south) -- (branch);
\draw[edge] (branch) -| (accept.north);
\draw[edge] (branch) -| (retain.north);
\node[labeltext,above] at (-1.30,-11.30) {是};
\node[labeltext,above] at (1.30,-11.30) {否};
\draw[stem] (accept.south) |- (outmerge);
\draw[stem] (retain.south) |- (outmerge);
\draw[edge] (outmerge) -- (output);
\end{tikzpicture}
\caption{MV-HGS-SP算法主流程}
\label{fig:algorithm-flow}
\end{figure}
